%% FullCourse_10Lecs -- Lecture 10, Topic 10.4: Case Study 3 --- ES Fellows Dataset
%% Source: ./archive_pre_split/Lec10_Case_Studies_v2.tex (section 4)
%% FullCourse-only content (not in MiniCourse).

%% Shared preamble for the FullCourse_10Lecs topic decks.
%%
%% Each topic .tex \input's this file, then defines its own \title, \date
%% override (if any), \graphicspath, and \begin{document}.
%%
%% Reconstructed verbatim from the inlined copy preserved in
%% Lec10_Case_Studies/Lec10_T8_Case6_DeepRamsey.tex (lines 23-190), which is
%% explicitly annotated there as "from ../shared_preamble.tex, copied verbatim".
%% Base preamble matches MiniCourse_8hr/Slides/shared_preamble.tex; this version
%% additionally provides \sectiondivider, the conceptbox/cautionbox/examplebox/
%% takeawaybox tcolorboxes, and \compacttables used across the split decks.

\documentclass[10pt,english,aspectratio=169]{beamer}

\usepackage{ae,aecompl}
\usepackage[T1]{fontenc}
\usepackage{textcomp}
\usepackage[utf8]{inputenc}
\usepackage{babel}
\usepackage{datetime2}

\setcounter{secnumdepth}{3}
\setcounter{tocdepth}{3}

\usepackage{amsmath, amssymb, amsthm, graphicx}
\usepackage{booktabs, multirow}
\usepackage{tikz}
\usepackage{pgfplots}
\pgfplotsset{compat=1.16}
\usepackage[most]{tcolorbox}
\tcbuselibrary{skins, breakable, theorems}
\usepackage{listings}
\usepackage{xcolor}

\lstset{
  basicstyle=\ttfamily\small,
  keywordstyle=\color{blue!80!black}\bfseries,
  commentstyle=\color{green!50!black}\itshape,
  stringstyle=\color{orange!80!black},
  backgroundcolor=\color{gray!8},
  frame=single,
  rulecolor=\color{gray!40},
  breaklines=true,
  showstringspaces=false,
  numbers=none,
}

\lstdefinestyle{pythonstyle}{
    language=Python,
    basicstyle=\ttfamily\footnotesize,
    keywordstyle=\color{blue!80!black}\bfseries,
    commentstyle=\color{green!50!black}\itshape,
    stringstyle=\color{orange!80!black},
    frame=single,
    breaklines=true,
    showstringspaces=false,
    numbers=left,
    numbersep=5pt,
    numberstyle=\tiny\color{gray},
    backgroundcolor=\color{gray!8},
    rulecolor=\color{gray!40}
}

\ifx\hypersetup\undefined
  \AtBeginDocument{%
    \hypersetup{bookmarks=false,breaklinks=true,pdfborder={0 0 0},colorlinks=true,
      linkcolor=DarkText,urlcolor=RedTitle}%
  }
\else
  \hypersetup{bookmarks=false,breaklinks=true,pdfborder={0 0 0},colorlinks=true,
    linkcolor=DarkText,urlcolor=RedTitle}%
\fi

\makeatletter
\newcommand\makebeamertitle{%
  \begin{frame}[plain]
    \vspace*{1.0cm}
    \centering
    {\usebeamerfont{title}\usebeamercolor[fg]{title}\inserttitle\par}
    \vspace{1.0cm}
    {\usebeamerfont{author}\usebeamercolor[fg]{author}\insertauthor\par}
    \vspace{0.2cm}
    {\usebeamerfont{date}\usebeamercolor[fg]{date}\insertdate\par}
  \end{frame}}%
\AtBeginDocument{%
  \let\origtableofcontents=\tableofcontents
  \def\tableofcontents{\@ifnextchar[{\origtableofcontents}{\gobbletableofcontents}}
  \def\gobbletableofcontents#1{\origtableofcontents}%
}

\usetheme{metropolis}
\usefonttheme{professionalfonts}

\definecolor{LightGray}{RGB}{230,230,230}
\definecolor{RedTitle}{RGB}{0,0,200}
\definecolor{VeryLightGray}{RGB}{245,245,245}
\definecolor{DarkText}{RGB}{0,0,0}

\setbeamercolor{background canvas}{bg=white}
\setbeamercolor{title}{fg=RedTitle}
\setbeamercolor{author}{fg=DarkText}
\setbeamercolor{date}{fg=DarkText}
\setbeamercolor{frametitle}{bg=VeryLightGray,fg=DarkText}
\setbeamercolor{normal text}{fg=DarkText}
\setbeamerfont{title}{size=\large}

\setbeamersize{text margin left=5mm, text margin right=5mm}
\addtobeamertemplate{frametitle}{}{\vspace{-0.5em}}
\newcommand{\reducevspace}{\vspace{-0.5cm}}

% Shared section divider and box styles for split topic decks.
\newcommand{\sectiondivider}[2]{%
\begin{frame}[plain,noframenumbering]
\begin{center}
\vfill
{\Large\textbf{#1}\par}
\vspace{0.35cm}
{\normalsize #2\par}
\vfill
\end{center}
\end{frame}
}

\newtcolorbox{conceptbox}[1][]{
  colback=blue!3,
  colframe=RedTitle!55,
  boxrule=0.35mm,
  arc=1mm,
  left=2mm,
  right=2mm,
  top=1mm,
  bottom=1mm,
  #1
}
\newtcolorbox{cautionbox}[1][]{
  colback=red!3,
  colframe=red!50!black,
  boxrule=0.35mm,
  arc=1mm,
  left=2mm,
  right=2mm,
  top=1mm,
  bottom=1mm,
  #1
}
\newtcolorbox{examplebox}[1][]{
  colback=green!3,
  colframe=green!45!black,
  boxrule=0.35mm,
  arc=1mm,
  left=2mm,
  right=2mm,
  top=1mm,
  bottom=1mm,
  #1
}
\newtcolorbox{takeawaybox}[1][]{
  colback=VeryLightGray,
  colframe=RedTitle!55,
  boxrule=0.35mm,
  arc=1mm,
  left=2mm,
  right=2mm,
  top=1mm,
  bottom=1mm,
  #1
}
\newcommand{\compacttables}{\renewcommand{\arraystretch}{1.15}}

\setbeamertemplate{navigation symbols}{}
\setbeamertemplate{footline}{%
  \hfill%
  \usebeamercolor[fg]{page number in head/foot}%
  \usebeamerfont{page number in head/foot}%
  \insertframenumber\,/\,\inserttotalframenumber\kern1em\vskip2pt%
}

\makeatother

\author{Zhigang Feng}
% "date with time": compile timestamp so each build is identifiable.
\date{\today\ \DTMcurrenttime}


% Suppress redundant auto-generated metropolis section page; \sectiondivider frames are the dividers we keep.
\metroset{sectionpage=none}

\graphicspath{{./pic/}{./figures/}{./Exercise10_ES_Fellows/plots/}{../../MiniCourse_8hr/Slides/Module4_Agentic_AI/pic/}{../}}

\usetikzlibrary{positioning,arrows.meta,shapes,calc,fit,backgrounds}

\lstdefinestyle{markdownstyle}{
    basicstyle=\ttfamily\footnotesize,
    frame=single,
    rulecolor=\color{gray!40},
    backgroundcolor=\color{gray!8},
    breaklines=true,
    showstringspaces=false,
    numbers=none,
}

\title[AI for Econ Research]{AI for Economic Research: Dynamic Models, Language, and Agents\\
Lecture 10.4: Case Study 3 --- ES Fellows Dataset}

\begin{document}

\makebeamertitle

%=============================================================================
\begin{frame}{Topic 10.4 Agenda}
\begin{enumerate}
  \item \textbf{Case Study 3} --- ES Fellows dataset construction with agentic AI
\end{enumerate}
\end{frame}

%=============================================================================
\section{Case Study 3: ES Fellows Dataset}
%===========================================================

%===========================================================

\sectiondivider{Case Study 3}{A provenance-aware public-web dataset construction workflow}

\begin{frame}{Why This Third Case Is Different}
\begin{itemize}
    \item the source material is public and semi-structured
    \item metadata quality is uneven across people and institutions
    \item missingness is part of the result, not a nuisance to hide
    \item the workflow must control sources, provenance, and coding rules
\end{itemize}

\vspace{0.3cm}
\textbf{This is not a web-scraping stunt.} It is a data-construction problem with explicit source restrictions.
\end{frame}

\begin{frame}{Research Question for the Fellows Dataset}
\begin{tcolorbox}[colback=blue!3, colframe=RedTitle!60]
\textbf{Question:} What can we learn about career stage, election timing, and field composition among current Econometric Society fellows using only official and institutional sources?
\end{tcolorbox}

\vspace{0.25cm}
\begin{itemize}
    \item it is a useful dataset about the economics profession itself
    \item it exposes a new research pattern: roster seed + metadata enrichment + coverage reporting
    \item it forces us to separate current biological age, years since PhD, and age at fellow election
    \item it forces students to confront missingness, source quality, and coding judgment
\end{itemize}
\end{frame}

\begin{frame}{Why a Naive Web-Chat Workflow Still Fails Here}
\begin{enumerate}
    \item \textbf{Roster control is weak.} You need a canonical seed list and fixed batches.
    \item \textbf{Source discipline is weak.} You need a rule that excludes Wikipedia and open-web guesswork.
    \item \textbf{State tracking is weak.} You need persistent files recording what was found and what stayed missing.
    \item \textbf{Auditability is weak.} You need field-level source URLs, not just narrative answers.
\end{enumerate}

\vspace{0.25cm}
\textbf{Precise contrast:} modern web chat can browse and summarize, but this assignment needs a project harness with files, batching, and provenance rules.
\end{frame}

\begin{frame}{Official Seed Roster as of April 14, 2026}
\begin{itemize}
    \item canonical source: the official Econometric Society current fellows page
    \item parsed roster size: 882 fellows
    \item official CV links on the seed page: 64
    \item usable election years in the current page: 880; two entries parse as \texttt{0000}
    \item batching rule: 18 batches of 50 fellows
    \item seed fields: name, affiliation, election year, profile URL, CV URL
\end{itemize}

\vspace{0.25cm}
\textbf{Interpretation:}
\begin{itemize}
    \item the roster is large enough that the workflow must be batched
    \item CV coverage is sparse enough that missingness will remain visible
    \item the current roster is not the same thing as a confirmed-living roster
\end{itemize}
\end{frame}

\begin{frame}{Batch Workflow and Project Structure}
\begin{columns}[T]
\begin{column}{0.48\textwidth}
\textbf{Project files}
\begin{itemize}
    \item seed roster CSV
    \item batch assignment file
    \item normalized batch CSV
    \item analysis script
    \item plots and memo
\end{itemize}
\end{column}
\begin{column}{0.48\textwidth}
\textbf{Workflow}
\begin{enumerate}
    \item parse the official roster
    \item assign a batch of 50
    \item fill official/institutional metadata
    \item run analysis
    \item write findings and limitations
\end{enumerate}
\end{column}
\end{columns}

\vspace{0.3cm}
\textbf{This is the public-web analog of `CLAUDE.md`:} the project files hold the rules and state across sessions.
\end{frame}

\begin{frame}{Schema and Provenance Rules}
\small
\begin{tabular}{p{3.4cm}p{8.4cm}}
\toprule
\textbf{Field} & \textbf{Rule} \\
\midrule
\texttt{name}, \texttt{affiliation}, \texttt{election\_year} & seed roster only \\
\texttt{bio\_url}, \texttt{cv\_url} & official or institutional only \\
\texttt{phd\_year}, \texttt{research\_fields\_raw} & record source URL for each field \\
normalized field bucket & map raw text into the lecture taxonomy \\
\texttt{coverage\_status} & say what remains missing \\
\texttt{notes} & document coding judgment or ambiguity \\
\bottomrule
\end{tabular}

\vspace{0.3cm}
\textbf{Credibility requires traceable provenance for each field.}
\end{frame}

\begin{frame}{Missingness, Gender, and Living-Status Caution}
\begin{itemize}
    \item \textbf{Birth date:} leave missing unless the official source states it
    \item \textbf{PhD year:} use only if the source states the degree year directly
    \item \textbf{Gender:} code only from explicit pronouns or explicit gender language
    \item \textbf{No inference:} not from names, photos, or country
    \item \textbf{Living status:} do not infer from old election year alone; for lecture purposes, use a 90-year cap only as a conservative sensitivity rule
\end{itemize}

\vspace{0.3cm}
\textbf{This is deliberate.} The assignment is designed to teach credible coding rules, not aggressive completion.
\end{frame}

\begin{frame}{Instructor Full Scan: What the Agentic Workflow Added}
\begin{columns}[T]
\begin{column}{0.44\textwidth}
\small
\begin{tabular}{lc}
\toprule
\textbf{Metric} & \textbf{Official $\rightarrow$ Web} \\
\midrule
Local CV docs & 62 $\rightarrow$ 316 \\
Birth year & 11 $\rightarrow$ 46 \\
PhD year & 47 $\rightarrow$ 183 \\
Field bucket & 15 $\rightarrow$ 73 \\
\bottomrule
\end{tabular}

\vspace{0.25cm}
\textbf{Main implication:}
\begin{itemize}
    \item the agent scaled discovery and parsing well
    \item PhD-year coverage improved materially
    \item birth-year coverage stayed thin
\end{itemize}
\end{column}
\begin{column}{0.54\textwidth}
\includegraphics[width=\textwidth]{full_web_coverage_dashboard.pdf}
\end{column}
\end{columns}
\end{frame}

\begin{frame}{The 90-Year Rule Does Not Solve the Main Problem}
\begin{columns}[T]
\begin{column}{0.55\textwidth}
\includegraphics[width=\textwidth]{full_web_age90_status.pdf}
\end{column}
\begin{column}{0.41\textwidth}
\small
\begin{itemize}
    \item observed age $\leq 90$: 44 fellows
    \item observed age $> 90$: 1 fellow
    \item birth year unknown: 836 fellows
\end{itemize}

\vspace{0.2cm}
\textbf{Takeaway:}
\begin{itemize}
    \item the problem is missing birth year, not a handful of extreme ages
    \item do not treat current biological age as the main variable
    \item use age at election or years since PhD instead
\end{itemize}
\end{column}
\end{columns}
\end{frame}

\begin{frame}{Career Stage: Use Years Since PhD and Age at Election}
\begin{columns}[T]
\begin{column}{0.5\textwidth}
\includegraphics[width=\textwidth]{full_web_years_since_phd.pdf}
\end{column}
\begin{column}{0.5\textwidth}
\includegraphics[width=\textwidth]{full_web_age_at_election.pdf}
\end{column}
\end{columns}

\vspace{0.15cm}
\textbf{Main descriptive findings:}
\begin{itemize}
    \item median years since PhD: 35
    \item median age at fellow election: 40
    \item median years from PhD to election: 17
    \item age at election is cleaner than current age because it does not depend on fellows being alive today
    \item this is exactly why agent-generated datasets still need variable-design judgment
\end{itemize}
\end{frame}

\begin{frame}{Election Counts Around Major Economics Events}
\begin{columns}[T]
\begin{column}{0.58\textwidth}
\includegraphics[width=\textwidth]{full_web_event_windows.pdf}
\end{column}
\begin{column}{0.38\textwidth}
\scriptsize
\textbf{3-year post-event windows}
\begin{itemize}
    \item Internet turn (2000): 41 $\rightarrow$ 45
    \item GFC (2008): 35 $\rightarrow$ 51
    \item COVID era (2020): 82 $\rightarrow$ 117
\end{itemize}

\vspace{0.2cm}
\textbf{Interpretation:}
\begin{itemize}
    \item there is no obvious post-dot-com collapse
    \item counts rise after 2008 and surge after 2020
    \item the 2020--2024 jump likely reflects Society-side policy or catch-up effects
    \item use this as descriptive timing evidence, not a causal event study
\end{itemize}
\end{column}
\end{columns}
\end{frame}

\begin{frame}{What Agentic AI Actually Did Here}
\begin{itemize}
    \item discovered and cached candidate homepages at scale, then kept the search state on disk
    \item parsed 527 institutional profile pages and downloaded 316 CV or CV-like documents
    \item normalized fields, tracked source URLs, and reran the analysis after rule changes
    \item surfaced a negative result clearly: even a large search pass does not rescue current biological age
\end{itemize}

\vspace{0.25cm}
\textbf{What stayed human:}
\begin{itemize}
    \item define source restrictions and coding rules
    \item separate current age, years since PhD, and age at election
    \item treat event patterns as descriptive rather than causal
    \item decide when missingness is too severe for a claim
\end{itemize}
\end{frame}

\begin{frame}{Exercise B: Full-Roster Fellows Assignment}
\begin{itemize}
    \item start from the official seed roster and take one batch of 50 fellows
    \item use only official and institutional sources
    \item produce a normalized CSV, four figures, and a short memo
    \item keep missing values visible and documented
    \item use age only if birth-year coverage reaches the lecture threshold; otherwise use years since PhD or age at election
    \item explain clearly what the agent automated and what judgment you supplied
\end{itemize}

\vspace{0.3cm}
\textbf{This case is not about scraping harder.} It is about designing a trustworthy public-web research pipeline.
\end{frame}


\end{document}
